% ESQUELETO DA DEFESA — estrutura apenas, sem conteúdo.
% Gerado a partir de PLANO_FLUXO_DEFESA.md §3, §8, §11 (2026-08-21).
% Nada de texto de slide entra aqui antes de 22/08 (plano validado com o orientador — ver §9).
% O showcase original do template (todos os comandos de exemplo: blocos, tabelas, código,
% algoritmos, equações, citações) foi preservado em template_showcase.tex — consulte-o para
% a sintaxe de qualquer recurso usado abaixo.

\documentclass{beamer}
\usepackage[net,green]{nesped}
\usepackage[brazilian]{babel}

% TODO: título exatamente como no folha de rosto entregue (src/) — fala em PT, slides em EN (§0).
\title[Título curto]{Multitask Learning for Point-of-Interest Classification and Prediction Tasks: The Role of the Check-in-Level Representation}
\subtitle{Defesa de Dissertação de Mestrado}
% TODO: confirmar grafia exata do nome como consta na folha de rosto.
\author{Vitor Hugo}
\institute[UFV]{Universidade Federal de Viçosa --- PPGCC}
\date{28 de agosto de 2026}

\begin{document}

\titleframe{
    \titlelogo{img/logo-nesped.png}
    \titlelogo{img/logo-ufv.png}
}

% Divisor automático de seção (§11.3). NUNCA \autotocframe{} com chaves vazias por acidente
% nem passar as opções sem colchetes -- o argumento vai para \tocframe[...], que é como
% \autotocframe repassa internamente (corrigido 2026-08-21, ver nesped.sty).
% subsectionstyle=hide, não show/show/shaded: com a profundidade de subseções deste plano
% (até 6 por seção), mostrar todas elas estoura o slide (validado 2026-08-21 -- 6 Overfull
% \vbox, um por seção, com show/show/shaded; 0 com hide). Ver relato ao agente ingred-14.
\autotocframe{sectionstyle=show/shaded, subsectionstyle=hide}

% ============================================================
% SEÇÃO 1 · Introdução — 4 min (PLANO §3, Seção 1)
% ============================================================
\section{Introdução}

\subsection{Motivação}
\begin{frame}{TODO 1.1 --- Motivação}
    % Regularidade da mobilidade (Song et al. ~93%, com a ressalva de não-teto na mesma frase)
    % + as aplicações que dependem de categoria e região.
\end{frame}

\subsection{Pergunta e veredito}
\begin{frame}{TODO 1.2 --- Pergunta de pesquisa e veredito}
    % A pergunta de pesquisa, literal (§1.2 do texto), e o veredito em uma linha,
    % em linguagem de lei (ver PLANO §5.1 --- NUNCA reformular).
\end{frame}

\subsection{Escopo}
\begin{frame}{TODO 1.3 --- Escopo}
    % Def. 2.7/2.8 dentro; Def. 2.9 (next-place) definida para ser EXCLUÍDA;
    % sete categorias; região = setor censitário / mahalle.
\end{frame}

\subsection{Roteiro}
\begin{frame}{TODO 1.4 --- Arco e roteiro}
    % O arco em uma linha + os três artigos/veículos + autoria (ver LO-11, PLANO §9.4).
\end{frame}

% TODO transição de saída do Ato I (PLANO §2, ato I) --- falada sobre o \autotocframe
% automático da próxima seção, ou \specialframe próprio se merecer tela cheia.

% ============================================================
% SEÇÃO 2 · Fundamentos compartilhados — 6 min (PLANO §3, Seção 2)
% Motor de de-duplicação do deck inteiro — nada disto é reaberto depois (PLANO §4).
% ============================================================
\section{Fundamentos}

\subsection{Linhagem}
\begin{frame}{TODO 2.1 --- Linhagem e representação}
    % Tabela 1 (DGI -> HGI -> Check2HGI; MTLnet -> ST-MTLNet -> modelo conjunto)
    % + contraste do título: vetor por lugar (Def. 2.3) x vetor por visita (Def. 2.4).
\end{frame}

\subsection{Vocabulário de MTL}
\begin{frame}{TODO 2.2 --- MTL, compartilhamento, transferência negativa}
    % Compartilhamento rígido (Def. 2.10), transferência negativa (Def. 2.12) como risco
    % nomeado, critério pré-registrado do balanceador (Nash-MTL etc. citados aqui, não depois).
\end{frame}

\subsection{Base de evidência}
\begin{frame}{TODO 2.3 --- Datasets}
    % Uma tabela, seis conjuntos (AL/AZ/FL/TX/CA + Istanbul), ordenada por nº de regiões,
    % dizendo qual capítulo usou quais.
\end{frame}

\subsection{Métricas}
\begin{frame}{TODO 2.4 --- Métricas}
    % macro-F1 (Food ~1/3, perda não reponderada), Acc@10 + desconto OOD, escore de seleção conjunta.
\end{frame}

\subsection{Protocolo}
\begin{frame}{TODO 2.5 --- Protocolo e regras de decisão}
    % Estratificado por amostra (Caps. 3/4) x disjunto por usuário (Cap. 5), no MESMO slide.
    % A lei dos verbos (§2.4.4 do texto).
\end{frame}

\subsection{Trabalhos relacionados}
\begin{frame}{TODO 2.6 --- Trabalhos relacionados}
    % Um slide, organizado pelo eixo que importa: categoria/região como MEIO para o
    % próximo lugar x como FIM. Nunca reaberto depois.
\end{frame}

% ============================================================
% SEÇÃO 3 · Capítulo 3 --- CBIC 2025: o nulo previsto — 4 min (PLANO §3, Seção 3)
% ============================================================
\section{CBIC 2025}

\subsection{Par de tarefas}
\begin{frame}{TODO 3.1 --- O par de tarefas}
    % Classificação estática (Def. 2.6) + próxima categoria --- um estático, um sequencial.
\end{frame}

\subsection{MTLnet}
\begin{frame}{TODO 3.2 --- MTLnet (Figura 1)}
    % Único desenho completo de arquitetura do deck: encoders por tarefa -> FiLM (uma cláusula)
    % -> tronco residual compartilhado -> duas cabeças.
\end{frame}

\subsection{Setup}
\begin{frame}{TODO 3.3 --- Setup}
    % Florida, sete categorias, 5 folds --- autodeclaração de protocolo dita aqui, não escondida.
\end{frame}

\subsection{Resultado}
\begin{frame}{TODO 3.4 --- O nulo}
    % Nas palavras do próprio capítulo: "largamente comparáveis, sem vantagem clara ou
    % consistente" + custo de treino.
\end{frame}

\subsection{Bifurcação}
\begin{frame}{TODO 3.5 --- Três hipóteses}
    % Dissimilaridade de tarefas / insuficiência de representação / rigidez de topologia.
\end{frame}

% TODO transição interna Cap.3 -> Cap.4 (PLANO §2, Ato II) --- congelar a arquitetura,
% mover só a entrada. Slide próprio ou falada sobre o \autotocframe da Seção 4.

% ============================================================
% SEÇÃO 4 · Capítulo 4 --- CoUrb 2026: o diagnóstico controlado — 4 min (PLANO §3, Seção 4)
% ============================================================
\section{CoUrb 2026}

\subsection{Pergunta herdada}
\begin{frame}{TODO 4.1 --- Arquitetura ou representação?}
    % MTLnet CONGELADO, só a entrada muda.
\end{frame}

\subsection{Ressalva da tarefa estática}
\begin{frame}{TODO 4.2 --- Ressalva antes do número}
    % A tarefa estática lê o próprio rótulo (venue-type -> 7 categorias 1:1) -> o ganho de
    % 20-22 pontos NÃO diz nada sobre a tarefa sequencial. Ressalva SEMPRE antes do número.
\end{frame}

\subsection{Resultado diagnóstico}
\begin{frame}{TODO 4.3 --- A tarefa sequencial}
    % O resultado diagnóstico de fato: alvo nunca está na entrada.
\end{frame}

\subsection{Bordas honestas}
\begin{frame}{TODO 4.4 --- Limites oferecidos}
    % Travel (topologia de grafo vence em movimento esparso de longa distância) +
    % comparação NÃO pareada em largura (192 x 64 dims).
\end{frame}

\subsection{Entrega}
\begin{frame}{TODO 4.5 --- A representação é o gargalo}
    % Frase de entrega, literal do Cap. 6: arquitetura fixa, entrada mudou, resultado moveu.
\end{frame}

% TODO transição de saída do Ato II (PLANO §2) --- o vazamento é o pivô, dito em voz alta
% aqui (enquadramento: a verificação funcionando, não uma confissão). Provável \specialframe.

% ============================================================
% SEÇÃO 5 · Capítulo 5 --- MobiWac 2026: a resolução — 20 min (PLANO §3, Seção 5)
% Quase metade da fala --- padrão curto-curto-longo. NÃO comprimir abaixo de ~18 min.
% ============================================================
\section{MobiWac 2026}

\subsection{O que muda}
\begin{frame}{TODO 5.1 --- O que muda e por quê (1 min)}
    % Representação (lugar -> check-in), topologia (rígida -> atenção cruzada),
    % protocolo (por amostra -> disjunto por usuário) --- como CONSEQUÊNCIAS do diagnóstico.
\end{frame}

\subsection{A representação}
\begin{frame}{TODO 5.2a --- Check2HGI (Figura 4, 5 min)}
    % Da sequência de check-ins às duas predições; o que um vetor por visita carrega;
    % a aresta só-para-frente (retomada do vazamento, uma única vez).
\end{frame}
\begin{frame}{TODO 5.2b --- Separabilidade (Figura 6)}
\end{frame}
\begin{frame}{TODO 5.2c --- Tabela 9}
    % Nas palavras da própria tabela (PLANO §5.2) --- não "bate nos seis", separa em CINCO.
\end{frame}

\subsection{A arquitetura}
\begin{frame}{TODO 5.3 --- Compartilhamento por troca (Figura 5, 3 min)}
    % Compartilhamento POR TROCA, não por camadas possuídas; caminho espacial privado da região.
    % A posição do tronco (PLANO §5.3) --- formulação simétrica, nunca "provavelmente contribuiu".
\end{frame}

\subsection{O protocolo}
\begin{frame}{TODO 5.4 --- Protocolo final (3 min)}
    % CV 5-fold disjunta por usuário, sementes {0,1,7,100}, 20 modelos ajustados,
    % teste pareado sobre as QUATRO médias por semente --- e o desvio declarado do plano.
\end{frame}

\subsection{O veredito}
\begin{frame}{TODO 5.5a --- Tabela 10 (5,5 min)}
    % Linguagem de lei (PLANO §5.1). "Supera" SÓ nas três células. Nunca "empata"/"ties"/
    % "em todos". As quatro células dentro da margem são déficits, não empates.
\end{frame}
\begin{frame}{TODO 5.5b --- Figura 7}
    % Cláusula obrigatória: ganhos de região em TX/CA são resultados secundários,
    % fora do plano registrado --- nunca creditados a transferência entre tarefas.
\end{frame}

\subsection{O trade}
\begin{frame}{TODO 5.6 --- Trade medido e limites (1,5 min)}
    % Os quatro limites declarados.
\end{frame}

% TODO transição de saída do Ato III (PLANO §2) --- "o que os três estudos, juntos,
% estabelecem -- e o que não estabelecem?"

% ============================================================
% SEÇÃO 6 · Conclusão Geral — 5 min (PLANO §3, Seção 6)
% ============================================================
\section{Conclusão}

\subsection{Resposta condicional}
\begin{frame}{TODO 6.1 --- A resposta condicional}
    % MTL ajudou SOB ESTE DESENHO E ESTE PROTOCOLO --- e o que isso não autoriza.
\end{frame}

\subsection{Contribuição}
\begin{frame}{TODO 6.2 --- A contribuição una}
    % Redação IDÊNTICA à do slide de abertura (1.2/1.4). Metade prática (um modelo,
    % uma passagem) + metade científica (as condições). Ganho operacional, não computacional.
\end{frame}

\subsection{Limitações e futuro}
\begin{frame}{TODO 6.3 --- Limitações x trabalhos futuros}
    % Um slide, duas colunas: seis limitações amarradas 1:1 aos seis trabalhos futuros.
    % Oferecidas antes de pedidas.
\end{frame}

\subsection{Fecho}
{
\specialframe
\begin{frame}{TODO 6.4 --- Fecho}
    % Retomada das aplicações + agradecimentos + linha do repositório.
\end{frame}
}

% ============================================================
% SÉRIE B --- trilha de reserva (PLANO §6). Fora da barra e da contagem do deck principal.
% Nenhum \section daqui em diante: isso tiraria a série B de "fora da contagem".
% ============================================================
\miniframesoff

\begin{frame}[noframenumbering]{B0 --- Índice}
    \hypertarget{b0}{}
    % TODO: um item por pergunta coberta, cada um \hyperlink{bXX}{...} para o slide B
    % correspondente. Cobertura obrigatória 1:1 com ARGUICAO.md (ver PLANO §6).
    \begin{itemize}
        \item \hyperlink{b1}{B1 --- Veredito e estatística (Q1, Q2, Q4, §7.1)}
        % \item \hyperlink{b2}{B2 --- Protocolo e vazamento (Q3, Q12, Q19)}
        % \item \hyperlink{b3}{B3 --- Pós-submissão e retratações (P1, Q13, Q14)}
        % \item \hyperlink{b4}{B4 --- Capítulos 3--4 (Q16--Q20)}
        % \item \hyperlink{b5}{B5 --- Não foi medido (U1--U8)}
        % \item \hyperlink{b6}{B6 --- Documento e escopo (errata do Resumo, Q15)}
    \end{itemize}
\end{frame}

\begin{frame}[noframenumbering]{B1 --- TODO: título = a pergunta, em português}
    \hypertarget{b1}{}
    % TODO: rodapé de proveniência obrigatório (página do volume principal / suplemento /
    % caminho de arquivo / "pós-submissão --- não consta em nenhum dos dois volumes").
    % Números copiados de célula impressa ou de ladder_recompute.json, nunca re-derivados.

    \hyperlink{b0}{\beamerreturnbutton{Voltar ao índice}}
\end{frame}

% TODO: replicar o padrão do slide B1 para B2--B6 (PLANO §6, tabela de famílias).

\end{document}
